% !TeX root = ../main.tex
\chapter{The CD4AI Pattern}
\label{chap:cd4ai}

CD4AI is a continuous delivery pattern that addresses the following problem:
\enquote{How can teams turn failures observed during the use of AI-enabled
systems into regression tests that support subsequent changes?}
This chapter presents the pattern and its feedback loop.
Chapter~\ref{chap:case-studies} examines how its stages appear in selected
systems and documented workflows.

\section{Context, Problem, and Intent}
\label{sec:cd4ai-context}

% Context

Foundation models are trained on broad data at scale and can be adapted to
a range of downstream tasks~\cite{Bommasani:FoundationModels:2021}.
Applications built around language foundation models can change their behavior
through modifications to prompts, retrieved context, tools, and orchestration,
without necessarily retraining the underlying
model~\cite{Schluntz:BuildingEffectiveAgents:2024}.
By avoiding a model-training cycle for such changes, application teams can
shorten the implementation cycle and introduce behavioral changes within their
ordinary software development workflow.

In this setting, evaluating AI-dependent behavior becomes part of the
application team's delivery responsibilities, even when the team does not
develop or train the underlying model. A change that is quick to implement
may still require substantial evaluation. The resulting challenge is to assess
its effects on established behavior while incorporating failures discovered
during use into future evaluations.

% Optional motivation: add a concrete observation from the startup experience,
% distinguishing personal observations from claims about development practices
% in general. Keep implementation details for the relevant case study.

% Problem

Initial evaluations cannot be assumed to cover every interaction that will
emerge during use. Teams therefore need a process for identifying candidate
failures, reviewing them, and turning confirmed defects into useful regression
cases. Collecting interactions alone does not establish this connection: the
cases must capture expected behavior and return to the evaluation suite used
to assess subsequent changes.

% Intent

CD4AI aims to support frequent delivery while preserving established behavior
by connecting testing, monitoring, and curation in a recurring feedback loop.
Although applications built around language foundation models motivate the
discussion here, the pattern is described in terms of these responsibilities,
rather than a particular model or evaluation framework. Its intent is to make
observed failures inform future validation, without guaranteeing that every
failure will be discovered or that a previously observed defect will never recur.

\section{Forces}
\label{sec:cd4ai-forces}

\noindent\emph{Working formulation: the following forces will be revisited
as the case studies are developed.}

Testability, robustness, and evolvability are desired qualities that can
reinforce one another. Testability helps teams assess whether a proposed change
preserves expected behavior before deployment. Robustness concerns the
application's ability to maintain acceptable behavior under varying conditions,
while evolvability concerns its ability to accommodate change. The evaluation
suite must also evolve to remain useful as the application changes.

These qualities need not oppose one another in every pairwise comparison.
The forces considered here instead express tensions between desired outcomes
and the constraints on achieving them. The pattern aims to address three
such tensions:

\begin{itemize}
  \item \textbf{Confidence in changes and delivery speed:} teams need evidence
    that changes preserve expected behavior, but running more extensive
    evaluations requires execution time and maintenance effort. The validation
    process must fit the time and resources available for frequent delivery.
  \item \textbf{Failure discovery and review capacity:} teams need to discover
    failures beyond those covered by existing tests, but collecting more
    candidate interactions increases the curation workload. Broad discovery
    must be balanced with the capacity to review candidates and prepare useful
    regression cases.
  \item \textbf{Behavior preservation and adaptation:} historical regression
    cases help check established behavior, but changing requirements can make
    some expectations outdated. Teams must distinguish unintended regressions
    from intended changes and maintain the suite accordingly.
\end{itemize}

Testing, monitoring, and curation work together to address these tensions.
Testing supplies evidence about proposed changes, monitoring supplies
interactions for investigation, and curation determines which interactions
should become regression cases and what behavior those cases should check.
The cycle does not remove the costs or trade-offs; its implementation must
establish how these responsibilities fit the team's delivery process.

% Writing notes:
% Keep the desired qualities distinct from the concrete tensions above.
% Do not claim that greater testability necessarily lowers robustness or
% evolvability. Better testing can support both.
% Validate this working formulation against Chapter 5: look for evidence of
% evaluation cost, candidate selection, review workload, and changes to expected
% behavior. Do not present these as findings before examining the evidence.
% Explain how the stages address the tensions without forcing a one-to-one
% correspondence between a stage and a force. Connected paragraphs are enough;
% pairwise-comparison subsections are not necessary.
% Pattern-writing guidance discussed during review:
% https://www.hillside.net/index.php/a-pattern-language-for-pattern-writing
% Add a formal citation if this guidance is used in the final argument.

\section{Overview of the Feedback Loop}
\label{sec:cd4ai-overview}

CD4AI proposes three stages to solve the problem and balance the forces:
testing, monitoring, and curation.

% Here we put the drawing


\begin{itemize}


  \item Introduce the three stages and their responsibilities before detailing
    individual techniques.
  \item Add a diagram: testing informs deployment; monitoring supplies candidate
    interactions; curation prepares regression cases that return to testing.
  \item Identify the information passed between stages and explain why collecting
    logs or running evaluations alone does not establish the complete cycle.
\end{itemize}

\section{Testing and Evaluation}
\label{sec:cd4ai-testing}

\subsection{Test Cases and Expected Behavior}
\begin{itemize}
  \item Explain how cases prepared for expected behavior complement regression
    cases derived from previously observed failures.
  \item Describe the input, relevant context, expected behavior, and evaluation
    criterion needed to make a retained case useful.
\end{itemize}

\subsection{Evaluation Criteria and Variability}
\begin{itemize}
  \item Clarify how tests and evaluations are used in this chapter, drawing on
    the terminology introduced in Chapter~\ref{chap:background}.
  \item Discuss possible checks and evaluators, including their limitations when
    several outputs are acceptable or responses vary between runs.
\end{itemize}

\subsection{Integration with the Delivery Pipeline}
\begin{itemize}
  \item Explain when evaluations run and how their results inform delivery.
  \item Separate the pattern's responsibility from implementation choices such
    as metrics, acceptance thresholds, and treatment of critical failures.
\end{itemize}

\section{Monitoring}
\label{sec:cd4ai-monitoring}

\subsection{Observing Interactions}
\begin{itemize}
  \item Describe what information can help reconstruct behavior during use:
    inputs, outputs, retrieved context, tool calls, and user feedback.
  \item Explain how the available observations constrain subsequent review.
\end{itemize}

\subsection{Selecting Candidates for Review}
\begin{itemize}
  \item Describe candidate signals, such as negative feedback, unexpected tool
    behavior, and guardrail violations, alongside sampling of interactions.
  \item Explain the trade-off between broad discovery and review effort, and
    distinguish a suspected failure from a confirmed defect.
\end{itemize}

\section{Curation}
\label{sec:cd4ai-curation}

\subsection{Reviewing Candidate Interactions}
\begin{itemize}
  \item Explain how human review, AI-assisted review, or a combination can assess
    selected interactions and distinguish defects from noise.
  \item Describe the role of expected behavior and review rationales; discuss
    uncertain judgments and evaluator limitations.
\end{itemize}

\subsection{Preparing Regression Cases}
\begin{itemize}
  \item Explain how a confirmed defect becomes a reusable input or scenario with
    an expected behavior and an evaluation criterion.
  \item Discuss which context must be retained and how cases are selected and
    maintained as the suite grows.
\end{itemize}

\subsection{Returning Cases to Testing}
\begin{itemize}
  \item Make the handoff explicit: recording a review does not by itself add a
    case to the regression suite.
  \item Describe manual and automated ways to retain the case, evaluate a fix,
    and check subsequent changes.
\end{itemize}

\section{Dynamics: Following a Failure through the Cycle}
\label{sec:cd4ai-dynamics}

\begin{itemize}
  \item Develop one clearly labeled illustrative example that follows an
    interaction from observation to review, a regression case, and evaluation
    of a subsequent change.
  \item Show the information and decisions at each transition, including where
    human intervention may be needed.
  \item Keep the example independent of a specific platform; reserve the
    evidence from concrete cases for Chapter~\ref{chap:case-studies}.
\end{itemize}

\section{Consequences and Limitations}
\label{sec:cd4ai-consequences}

\begin{itemize}
  \item Discuss the intended benefit of a regression suite informed by actual
    use, separating expected benefits from measured outcomes.
  \item Discuss the costs of observation, review, and maintaining and running
    the growing suite.
  \item Explain limits arising from unobserved interactions, missed signals,
    reviewer errors, and variable behavior. Avoid guarantees of complete
    coverage or prevention of every recurrence.
  \item Introduce the case studies as an examination of how these responsibilities
    and connections appear in different systems and workflows.
\end{itemize}
